\subsection{State of the art}
Earlier work on Magnus' conjecture was done by Conway and one of the authors in \cite{CD} using Riemann--Hilbert (RH) techniques. The weight function considered therein had the form
\begin{align}\label{wk}
 w_k(x)= \log \left(\frac{2k}{1-x}\right), \qquad x \in [-1,1), \quad k > 1.
\end{align}
The authors prove the conjecture for this special case corresponding to $\alpha = \beta = 0$, and also obtain $B = 0$ and $C = -\frac{3}{32}$, further suggesting that these constants do not depend on the behaviour of the weight function away from the logarithmic singularity.

\begin{Theorem}[{\cite{CD}}]\label{CDTheorem}
 The recurrence coefficients $\big\lbrace a_n^{(k)}\big\rbrace_{n=0}^\infty$, $\big\lbrace b_n^{(k)} \big\rbrace_{n=0}^\infty$ of the orthogonal polynomials with orthogonality measure $w_k(x){\rm d}x$, with $w_k$ given in \eqref{wk}, satisfy
\begin{align*}
 a_n^{(k)} = -\frac{3}{16n^2\log^2 n} + O\left( \frac{1}{n^2\log^3 n}\right) \qquad \text{as} \quad n \to \infty
\end{align*}
and
\begin{align*}
 b_n^{(k)} = \frac{1}{2}+\frac{1}{16n^2}-\frac{3}{32n^2\log^2 n} + O\left(\frac{1}{n^2\log^3 n}\right) \qquad \text{as} \quad n \to \infty.
\end{align*}
\end{Theorem}
Note that for $k > 1$, the weight function $w_k(x)$ has a positive finite value at $x = -1$, while $\lim_{k\to 1+} w_k(x) = w(x)$ for $x \in [-1,1)$. Hence, we see different $1/n^2$-terms in Theorem \ref{MainTheorem} compared to Theorem \ref{CDTheorem}, consistent with the conjectural \eqref{MagA} and \eqref{MagB}. Also observe that terms of order $1/\big(n^2 \log^2 n\big)$ are identical in both theorems, in particular they do not depend on~$k$. These terms can be interpreted as contributions from the logarithmic singularity at $x = +1$, which does not depend on $k \geq 1$.

For $k < 1$, the weight function $w_k$ would have a simple zero in the interior of $(-1,1)$ and the uniqueness of the corresponding polynomials is no longer guaranteed. We will not deal with this case in the present paper.

The main difficulty encountered in \cite{CD} was the lack of a known parametrix in the vicinity of the logarithmic singularity, a key ingredient in the usual nonlinear steepest descent analysis. Hence, the authors relied instead on a technically involved comparison to the Legendre problem with weight function $w_{Leg}(x) \equiv 1$, $x \in (-1,1)$. Surprisingly, their argument could not be generalized in an obvious way to the weight function $w = \lim_{k\to 1} w_k$, due to the appearance of a~simple zero of $w$ at $x = -1$. While, an analogous comparison to the Jacobi problem with weight function~$w_{Jac}(x) = 1+x$ (or something similar) seems suggestive, significant challenges remain due to the presence of the simple zero. In particular, the crucial Theorem 4.7 in \cite{CD} requires a different proof in this case: now one needs to control the behaviour of the Cauchy operator acting on spaces with Muckenhoupt weights.

Orthogonal polynomials with logarithmic weight functions, have applications to both pure mathematics and physics. In particular, apart from logarithmic singularities, these weight functions also tend to have zeros (see \cite[Section~8]{Mag18} and \cite{VA}). Such applications motivate us in the present paper to extend the results in~\cite{CD} to the weight function $w$ in~\eqref{wlog}, having both a~logarithmic singularity and a simple zero.

\begin{Remark}
 One might think, a priori, that the vanishing of a weight $w(x)$ at a point should not give rise to serious technical difficulties. Naively, it would appear that only singularities in the weight, and not zeros, should present obstacles. In this regard, we recall the hope and the prophecy of Lenard's 1972 paper~\cite{Lenard}:
 \begin{quote}
 ``It is the author’s hope that a
rigorous analysis will someday carry the results to the point where the true role of the zeros of the
generating function will be understood. When that day comes a capstone will have been put on a
beautiful edifice to whose construction many contributed and whose foundations lie in the studies
of Gabor Szeg\H{o} half a century ago''.
\end{quote}
\end{Remark}

\subsection[Relation of the present work to cite\{CD\}]{Relation of the present work to \cite{CD}}

Significant parts of the analysis performed in the present paper are based on the analysis introduced in \cite{CD}. Hence, we will repeatedly refer to that paper for proofs of certain statements. This is justified by the fact that the majority of estimates found in \cite{CD} do not depend on the distinction $k > 1$ and $k = 1$ in \eqref{wk}. Thus, the proofs of many of the results will also hold for the weight function \eqref{wlog} that we are interested in. There are however certain propositions which have their analogs in~\cite{CD}, but still deserve a separate proof due to some minor differences. These are Proposition~\ref{Proprr} which is the analog of Proposition~2.5 in~\cite{CD}, and Proposition~\ref{IntegralEst2} which is the analog of Propositions~5.3 and~5.4 in~\cite{CD}. In the case of Proposition~\ref{Proprr}, it is necessary to prove a slightly more general result than Proposition~2.5 in \cite{CD}. Meanwhile, Proposition~6.5 contains an application of Proposition~\ref{Proprr} in its more general form, and uses different RH solutions than Propositions~5.3 and 5.4 in \cite{CD}. Both results are proven in Appendix~\ref{appendixA}.

